\chapter{Семейство моделей GigaChat}\label{ch:gigachat}

Генеративные большие языковые модели стали важнейшим инструментом для исследований и приложений в области обработки естественного языка на различных языках. Однако разработка фундаментальных моделей, специально адаптированных для русского языка, остаётся ограниченной, прежде всего из-за значительных требований к вычислительным ресурсам. В настоящей главе представлено семейство русскоязычных LLM GigaChat, включающее модели различных масштабов: базовые предобученные модели и версии, дообученные на инструкциях. Приводится детальное описание архитектуры моделей, процесса предобучения и экспериментов, определивших ключевые проектные решения. Результаты настоящей главы опубликованы в~\cite{2506.09440}.

Семейство GigaChat разработано коллективом авторов и в настоящей работе выступает методологической и экспериментальной основой: оно поставляет предобученные бэкбоны для модели эмбеддингов (Глава~\ref{ch:gigaembeddings}), реранкера (Глава~\ref{ch:reranking}) и классификатора безопасности (Глава~\ref{ch:safety}). Поэтому глава носит характер описания этой основы, а не самостоятельного выносимого на защиту результата (см. личный вклад автора во Введении).

%% =================================================================
\section{Архитектура модели}\label{sec:gc_architecture}
%% =================================================================

\subsection{Обзор семейства моделей}\label{subsec:gc_overview}

Семейство GigaChat является первой коллекцией фундаментальных и пост-обученных моделей, специально разработанных и предобученных с нуля для русского языка. Начальная версия семейства GigaChat использует архитектуру Mixture of Experts (MoE) и включает три модели, выпущенные в открытый доступ: базовую модель, версию, дообученную на инструкциях, и версию, выровненную с помощью Direct Preference Optimization (DPO)~\cite{rafailov2023dpo}. Продвинутые проприетарные модели --- Lite, Pro и MAX --- непрерывно обновляются и доступны через пользовательский API и Telegram-бот.
Семейство развивается: последующие поколения (GigaChat~2, GigaChat~3) расширяют линейку как по масштабу, так и по архитектуре, охватывая модели от сотен миллионов до десятков миллиардов параметров.
В настоящей работе в качестве бэкбонов используются модели различных размеров из этого семейства (Таблица~\ref{tab:gc_usage}): компактные декодерные модели GigaChat-0.6B и GigaChat-3B для поисковых компонентов (Главы~\ref{ch:gigaembeddings}--\ref{ch:reranking}) и MoE-модель GigaChat3.1-10B-A1.8B для классификатора безопасности (Глава~\ref{ch:safety}).

\begin{table}[H]
    \centering
    \caption{Модели семейства GigaChat, используемые в настоящей диссертации.}
    \label{tab:gc_usage}
    \begin{tabular}{p{0.34\textwidth}p{0.56\textwidth}}
        \toprule
        Модель & Где используется в диссертации \\
        \midrule
        GigaChat-A3B-base       & базовая MoE-модель семейства (настоящая глава) \\
        GigaChat-3B             & бэкбон GigaEmbeddings (Глава~\ref{ch:gigaembeddings}) и реранкера (Глава~\ref{ch:reranking}) \\
        GigaChat-0.6B           & бэкбон компактного реранкера (Глава~\ref{ch:reranking}) \\
        GigaChat3.1-10B-A1.8B   & классификатор безопасности (Глава~\ref{ch:safety}) \\
        \bottomrule
    \end{tabular}
\end{table}

\subsection{Архитектура MoE}\label{subsec:gc_moe}

Модель GigaChat-A3B-base использует архитектуру MoE с 20 миллиардами общих параметров, из которых приблизительно 3.3 миллиарда активируются при каждом прямом проходе (Таблица~\ref{tab:gc_arch}). В экспериментах на одних и тех же данных архитектура MoE продемонстрировала значительное повышение эффективности, включая двукратное увеличение скорости обучения и снижение задержки инференса на 40\% по сравнению с плотными моделями аналогичного масштаба, такими как LLaMA~3 8B.

\begin{table}[H]
    \centering
    \caption{Конфигурация архитектуры модели GigaChat-A3B-base.}
    \label{tab:gc_arch}
    \begin{tabular}{lc}
        \toprule
        Параметр & Значение \\
        \midrule
        Архитектура                    & MoE \\
        Общее число параметров         & 20 млрд \\
        Активируемых параметров        & $\approx$3.3 млрд \\
        Число слоёв                    & 28 \\
        Общие эксперты                 & 2 \\
        Маршрутизируемые эксперты      & 64 \\
        KV-головы                      & 8 \\
        Головы внимания                & 16 \\
        Длина контекста                & 131k \\
        \bottomrule
    \end{tabular}
\end{table}

Эффективность обеспечивается блочно-разреженными вычислениями с использованием оптимизированных ядер STK Triton (вместо Megablocks) и выборочным контрольным сохранением активаций, что снижает вычислительные требования обучения на 40\% по сравнению с плотной моделью сопоставимого масштаба при обработке 1 триллиона токенов. Данные оптимизации устраняют необходимость в параллелизме экспертов при сохранении производительности модели.

Архитектура заменяет стандартные MLP-блоки слоями MoE (за исключением первого слоя, использующего gated MLP из-за проблем распределения токенов). Каждый MoE-блок использует множество экспертов и ненормализованный маршрутизатор (router) для обеспечения специализации, следуя принципам DeepSeek MoE~\cite{dai2024deepseekmoe}. Промежуточная размерность расширена до 14\,336 (аналогично Mistral~7B~\cite{jiang2023mistral}) для увеличения ёмкости, а эксперты разделяются между слоями для повышения параметрической эффективности.

\subsection{Анализ специализации экспертов}\label{subsec:gc_experts}

Для исследования доменной специализации экспертов был проведён анализ поведения маршрутизатора на подмножестве датасета Pile~\cite{gao2020pile}. На основе активаций маршрутизатора были построены эмбеддинги: для каждого образца формировалась матрица размером $l \times e$, где $l$ --- число MoE-слоёв, $e$ --- число экспертов в слое. Значение $\text{emb}_{ij}$ вычислялось как число активаций эксперта $j$ в слое $i$, нормализованное длиной образца в токенах.

Кластеризация эмбеддингов с помощью UMAP и HDBSCAN показала, что образцы группируются по доменам (Рисунок~\ref{fig:gc_umap}), что свидетельствует о кодировании доменной информации решениями маршрутизатора. Были идентифицированы кластеры в доменах спорта, кулинарии, биологии и программирования.

\begin{figure}[H]
    \centering
    \includegraphics[width=\textwidth]{data/gigachat-family-efficient-russian-language-modeling-through-mixture-of-experts-architecture/figure-14-2.jpg}
    \caption{Двумерная проекция эмбеддингов маршрутизатора с помощью UMAP. Цвета соответствуют различным доменам датасета Pile.}
    \label{fig:gc_umap}
\end{figure}

На основе данного анализа были созданы доменно-специфичные эмбеддинги путём усреднения значений внутри кластеров. Показано, что данный метод позволяет управлять процессом генерации: например, использование эмбеддингов, связанных со спортивной тематикой, приводило к генерации текстов спортивной направленности.

%% =================================================================
\section{Токенизация}\label{sec:gc_tokenizer}
%% =================================================================

Для семейства GigaChat был разработан специализированный токенизатор, повышающий эффективность кодирования кириллических слов, языков программирования и LaTeX. Точность обработки кода улучшена за счёт включения распространённых ключевых слов и поддержки пробелов, табуляций и переносов строк. Высокочастотные термины из LaTeX и языков программирования включены в словарь для минимизации фрагментации.

Для обучения токенизатора применён алгоритм Byte-Pair Encoding (BBPE) от Hugging Face. В ходе экспериментов постепенно корректировалась пропорция текстов из различных доменов (русский, английский, другие языки, код), что привело к генерации более ста кандидатов-токенизаторов, из которых был выбран наилучший по среднему отношению символов на токен. Для обучения использовались открытые датасеты: FRW, RedPajama, StarCoder, а также данные из Common Crawl, Wikipedia и Stack Exchange.

Результаты сравнения токенизаторов по среднему отношению символов на токен в различных доменах представлены в Таблице~\ref{tab:gc_tokenizer}. Токенизаторы GigaChat демонстрируют наилучшие средние показатели, превосходя токенизаторы GPT-4o, LLaMA~3, Qwen2 и других моделей.

\begin{table}[H]
    \centering
    \caption{Сравнение токенизаторов по отношению числа символов на токен. Токенизаторы \texttt{giga\_tokenizer} представляют варианты из экспериментов GigaChat.}
    \label{tab:gc_tokenizer}
    \resizebox{\textwidth}{!}{%
    \begin{tabular}{ccccccccc}
        \toprule
        \multirow{3}{*}{\textbf{Токенизатор}} & \multicolumn{3}{c}{\textbf{Код}} & \multirow{3}{*}{\textbf{ArXiv}} & \multicolumn{3}{c}{\textbf{Wiki}} & \multirow{3}{*}{\makecell{\textbf{Средний}\\\textbf{балл}}} \\
        \cmidrule(lr){2-4} \cmidrule(lr){6-8}
          & C & Java & C\# &  & Ru & Ar & En &  \\
        \midrule
        giga\_tokenizer\_1 & 3.57 & 4.15 & 4.62 & \textbf{3.61} & \underline{4.18} & \underline{3.34} & 4.47 & \textbf{3.99} \\
        giga\_tokenizer\_2 & 3.56 & 4.14 & 4.60 & \textbf{3.61} & 4.14 & 3.30 & 4.44 & \underline{3.97} \\
        gpt-4o            & \underline{3.74} & 4.43 & 4.88 & 3.39 & 3.40 & 3.07 & \textbf{4.68} & 3.94 \\
        giga\_tokenizer\_5 & 3.39 & 3.97 & 4.44 & \underline{3.54} & \textbf{4.20} & \textbf{3.50} & 4.43 & 3.92 \\
        giga\_tokenizer\_3 & 3.51 & 4.11 & 4.59 & \underline{3.54} & 4.04 & 3.25 & 4.35 & 3.91 \\
        giga\_tokenizer\_4 & 3.50 & 4.11 & 4.58 & 3.53 & 4.00 & 3.21 & 4.33 & 3.90 \\
        llama-3           & \textbf{3.75} & \underline{4.54} & \textbf{4.99} & 3.38 & 3.02 & 2.60 & 4.62 & 3.85 \\
        mistral-nemo      & 3.38 & 4.06 & 4.50 & 3.49 & 3.18 & 3.24 & 4.51 & 3.76 \\
        qwen2             & 3.69 & 4.52 & 4.95 & 3.31 & 2.70 & 2.56 & 4.50 & 3.75 \\
        gpt-4             & \underline{3.74} & \textbf{4.55} & \underline{4.98} & 3.38 & 2.04 & 1.44 & \underline{4.62} & 3.54 \\
        nemotron-4-256k   & 2.82 & 3.34 & 3.76 & 3.25 & 3.20 & 2.93 & 4.57 & 3.41 \\
        deepseek-coder-v2 & 2.95 & 3.51 & 3.92 & 3.35 & 2.39 & 1.11 & 4.42 & 3.10 \\
        deepseek-v2       & 2.95 & 3.51 & 3.92 & 3.35 & 2.39 & 1.11 & 4.42 & 3.10 \\
        mistral-large     & 2.75 & 3.26 & 3.64 & 3.14 & 2.46 & 1.13 & 4.04 & 2.92 \\
        \bottomrule
    \end{tabular}%
    }
\end{table}

%% =================================================================
\section{Данные для предобучения}\label{sec:gc_pretrain_data}
%% =================================================================

Предобучающий корпус сформирован из разнообразных текстовых источников с обеспечением баланса между лингвистическим богатством, доменно-специфичными знаниями и качеством данных. Статистика данных представлена в Таблице~\ref{tab:gc_pretrain_data}.

\begin{table}[H]
    \centering
    \caption{Статистика предобучающих данных по источникам.}
    \label{tab:gc_pretrain_data}
    \begin{tabular}{lcc}
        \toprule
        Источник данных & Уникальных токенов & Просмотрено токенов \\
        \midrule
        Веб                          & 4.4T & 5.6T \\
        Высококачественные источники & 630B & 1.3T \\
        Программный код              & 230B & 1.3T \\
        Синтетические данные         & 9B   & 81B  \\
        \bottomrule
    \end{tabular}
\end{table}

\textbf{Веб-данные.} Корпус высококачественных веб-данных сформирован на основе дампов Common Crawl за 2017--2023~гг. с использованием легковесного классификатора для извлечения многоязычных текстов на русском, английском, казахском, узбекском, португальском и арабском языках. Тексты дополнительно классифицированы с помощью LLM и специализированных моделей для идентификации образовательного и высокоинформативного контента, что позволило получить 4.4T уникальных токенов. Датасет преимущественно англоязычный (63.76\%) и русскоязычный (26.49\%).

\textbf{Высококачественные текстовые источники.} Включены книги и научные статьи открытого доступа, обработанные с помощью продвинутого оптического распознавания символов, а также научные и энциклопедические источники (arXiv, Wikipedia, PubMed), что добавило 630B токенов.

\textbf{Программный код.} Использован датасет StarCoder2~\cite{lozhkov2024starcoder2} совместно с курируемым набором открытого программного обеспечения. Модели машинного обучения применены для фильтрации низкокачественного кода, что позволило получить подмножество в 230B токенов.

\textbf{Синтетические данные.} Для математики и программирования построен конвейер, расширяющий исходные математические задачи путём многократного решения с фильтрацией через голосование большинством. Также созданы высококачественные синтетические задачи по программированию (сложные задачи на Python с документацией, объяснениями и проверками), диверсифицированные с использованием персон и липограмм.

Во всех источниках реализована точная дедупликация, а для англоязычных данных дополнительно применена MinHash-дедупликация~\cite{broder1997resemblance}.

%% =================================================================
\section{Процесс обучения}\label{sec:gc_training}
%% =================================================================

\subsection{Предобучение}\label{subsec:gc_pretrain}

Базовая модель обучена с использованием постоянного многошагового планировщика скорости обучения с разогревом. Планировщик включал период разогрева в 2000 батчей, после чего выполнялись четыре шага снижения скорости обучения на 30\%, 60\%, 90\% и 98\% от общей продолжительности обучения. На каждом шаге скорость обучения умножалась на $0.25^k$, где $k = 1, 2, 3, 4$ соответственно. Начальная скорость обучения составляла $10^{-4}$.

Процесс обучения использовал глобальный размер батча приблизительно 16 миллионов токенов (2048 последовательностей по 8192 токена), причём по мере обучения размер батча увеличивался. Основная фаза предобучения охватила приблизительно 9.5 триллиона токенов; общий объём обработанных токенов, включая финальную фазу отжига (annealing), составил приблизительно 10 триллионов.

После начального этапа обучения проведено расширение контекста в два этапа: сначала до 32K, затем до 128K токенов. Для улучшения работы с расширенным контекстом скорректирована база RoPE-эмбеддингов~\cite{su2024rope} с использованием подхода ABF~\cite{xiong2023effective}: 10K для начального контекста 8K, 300K для 32K и 1.4M для 128K. Расширение контекста потребовало около 1.8 триллиона дополнительных токенов.

Использована методология обучения со смешанной точностью (bfloat16 для большинства операций и fp32 для критических компонентов, таких как маршрутизатор). Для преодоления узких мест коммуникации в крупномасштабных распределённых средах с более чем 256 GPU размер батча увеличивался вместо добавления устройств с тем же объёмом работы.

\subsection{Дообучение на инструкциях}\label{subsec:gc_posttrain}

Данные для дообучения аннотированы профессиональными AI-тренерами, оценивавшими ответы по критериям следования инструкциям, учёта контекста, фактической точности и безопасности. Для генерации разнообразных диалогов в различных доменах создан проект аннотирования <<Dialog Creation>> на краудсорсинговой платформе Tagme.

Для улучшения способностей модели к запоминанию и извлечению информации применена генерация с извлечением информации (RAG), следуя экспериментам Grattafiori et al.~\cite{grattafiori2024llama3}. Данный подход генерирует доменно-специфичные обучающие данные из предобучающего корпуса, повышая контекстное понимание.

Пост-обучение модели GigaChat-A3B-instruct включает приблизительно 250 тысяч примеров. Распределение по доменам представлено в Таблице~\ref{tab:gc_sft_data}. Гиперпараметры дообучения представлены в Таблице~\ref{tab:gc_sft_hp}.

\begin{table}[H]
    \centering
    \caption{Распределение доменов инструкций в данных дообучения GigaChat-A3B-instruct.}
    \label{tab:gc_sft_data}
    \begin{tabular}{lc}
        \toprule
        Домен & Доля \\
        \midrule
        диалоги                   & 10\% \\
        длинный контекст (книги)  & 4\%  \\
        программный код           & 4\%  \\
        наука                     & 16\% \\
        общие знания о мире (веб) & 34\% \\
        переводы                  & 1\%  \\
        редактирование текста     & 12\% \\
        бизнес-специфика          & 3\%  \\
        вызовы функций / API      & 16\% \\
        \bottomrule
    \end{tabular}
\end{table}

\begin{table}[H]
    \centering
    \caption{Гиперпараметры дообучения моделей GigaChat.}
    \label{tab:gc_sft_hp}
    \resizebox{\textwidth}{!}{%
    \begin{tabular}{lllll}
        \toprule
        Модель & Оптимизатор & Планировщик & Параметры планировщика & Гиперпараметры \\
        \midrule
        GigaChat-A3B-instruct     & AdamW & постоянный & пользовательское снижение              & --- \\
        GigaChat-A3B-instruct~1.5 & AdamW & косинусный & разогрев: 200 шагов, макс. 7900 шагов  & betas (0.9, 0.95), eps: 1.0e-8 \\
        \bottomrule
    \end{tabular}%
    }
\end{table}

\subsection{Оптимизация прямых предпочтений}\label{subsec:gc_dpo}

При разработке GigaChat-A3B-instruct~1.5 были выявлены ключевые проблемы стандартного DPO: фокус на увеличении разрыва между хорошими и плохими ответами вместо улучшения точности, что приводит к галлюцинациям и нестабильности, а также игнорирование важности общих токенных префиксов.

Для решения этих проблем предложены модификации функции потерь DPO, включающие уникальные весовые коэффициенты, приоритизирующие улучшение качественных ответов над подавлением некачественных, с особым вниманием к общим префиксам. Добавлен нормализованный термин отрицательного логарифмического правдоподобия относительно референсной модели для стабилизации соотношений потерь.

%% =================================================================
\section{Экспериментальная оценка}\label{sec:gc_evaluation}
%% =================================================================

\subsection{Бенчмарки и методология}\label{subsec:gc_benchmarks}

Для оценки моделей используются общепринятые бенчмарки на английском и русском языках: математические способности (GSM8K~\cite{cobbe2021gsm8k}, MATH~\cite{hendrycks2021math}), программирование (HumanEval~\cite{chen2021humaneval}, MBPP), общие знания (MMLU EN~\cite{hendrycks2020mmlu}, MMLU RU, MMLU PRO~\cite{wang2024mmlupro}, RUBQ~\cite{korablinov2020rubq}, WINOGRANDE~\cite{sakaguchi2021winogrande}), кибербезопасность (CyberMetric~\cite{tihanyi2024cybermetric}) и следование инструкциям (IFEval~\cite{zhou2023ifeval}). Для русскоязычной оценки использован бенчмарк MERA~\cite{fenogenova2024mera}. Для всех тестов применён фреймворк LM Evaluation Harness.

\subsection{Результаты на англоязычных бенчмарках}\label{subsec:gc_english}

Результаты сравнения с открытыми пост-обученными LLM совместимых размеров представлены в Таблице~\ref{tab:gc_benchmarks}. Модели GigaChat-A3B-instruct и GigaChat-A3B-instruct~1.5 (3.3B активных параметров) демонстрируют сбалансированное соотношение масштаба и производительности по сравнению с моделями на 7--8B параметров (Qwen2.5~7B, Llama~3.1~8B, T-Lite).

Отмечаются разрывы в математических ($-14\%$ GSM8K, $-34\%$ MATH) и программистских ($-46\%$ MBPP, $-55\%$ HumanEval) задачах, отражающие параметрические ограничения, при превосходстве в рассуждениях ($+15\%$ RUBQ, $+12\%$ WINOGRANDE) и сохранении конкурентоспособности на MMLU ($-5\%$ до $-8\%$). Модели GigaChat~2~MAX и GigaChat~2~Pro демонстрируют лучшие результаты на всех бенчмарках, лишь незначительно уступая на CyberMetric ($-5\%$).

\begin{table}[H]
    \centering
    \caption{Сравнение моделей на русско- и англоязычных бенчмарках. Лучший результат в каждом столбце выделен жирным, лучший результат в том же размерном классе подчёркнут.}
    \label{tab:gc_benchmarks}
    \resizebox{\textwidth}{!}{%
    \begin{tabular}{l|l|lllll|ll}
        \toprule
        \textbf{Бенчмарк} & \textbf{Примеров} & \thead{GigaChat-\\A3B-instruct} & \thead{GigaChat-A3B\\-instruct 1.5} & \textbf{Qwen 2.5} & \textbf{T Lite} & \textbf{Llama 3.1} & \textbf{GigaChat2 Pro} & \textbf{GigaChat2 MAX} \\
        \midrule
        GSM8K & 5 & 0.764 & 0.774 & \underline{0.895} & 0.882 & 0.789 & 0.95 & \textbf{0.956}\\
        MATH & 4 & 0.462 & 0.393 & \underline{0.704} & 0.592 & 0.329 & 0.752 & \textbf{0.773}\\
        \midrule
        HumanEval & 0 & 0.329 & 0.378 & \underline{0.854} & 0.799 & 0.683 & \textbf{0.915} & 0.871\\
        MBPP & 0 & 0.385 & 0.441 & \underline{0.820} & 0.759 & 0.725 & 0.862 & \textbf{0.894} \\
        \midrule
        MMLU EN & 5 & 0.648 & 0.650 & \underline{0.710} & 0.718 & 0.682 & 0.821 & \textbf{0.86} \\
        MMLU RU & 5 & 0.598 & 0.600 & \underline{0.632} & 0.626 & 0.569 & 0.775 & \textbf{0.805} \\
        MMLU PRO EN & 5 & 0.348 & 0.357 & \underline{0.565} & 0.509 & 0.443 & 0.644 & \textbf{0.667} \\
        RUBQ & 0 & 0.675 & \underline{0.688} & 0.373 & 0.583 & 0.484 & 0.658 & \textbf{0.723} \\
        WINOGRANDE & 4 & 0.750 & \underline{0.762} & 0.636 & 0.670 & 0.624 & 0.796 & \textbf{0.832} \\
        CyberMetric & 0 & 0.798 & 0.791 & 0.787 & \underline{0.883} & 0.796 & \textbf{0.84} & 0.832 \\
        \midrule
        IFEval & 0 & 0.411 & 0.433 & \underline{0.819} & 0.730 & 0.812 & 0.837 & \textbf{0.899} \\
        \bottomrule
    \end{tabular}%
    }
\end{table}

\subsection{Результаты на русскоязычных бенчмарках}\label{subsec:gc_russian}

Результаты оценки на бенчмарке MERA представлены в Таблице~\ref{tab:gc_mera}. Модели GigaChat~2~MAX, GigaChat~2~Pro и GigaChat~2 достигают результатов близких к state-of-the-art ($\pm 2{-}7\%$) и доминируют в специализированных задачах: ruModAr ($+4\%$ до $+29\%$) и USE ($+7\%$ до $+33\%$), демонстрируя сильные стороны в логике и сложном понимании текста. Разрешение кореференций (RWSD: $-7\%$ до $-18\%$) и продвинутые рассуждения (MaMuRAMu: $-3\%$ до $-4\%$) показывают потенциал для роста, при этом производительность остаётся конкурентоспособной как по сравнению с фронтирными моделями (GPT-4o), так и с альтернативами среднего класса.

Модель GigaChat-A3B-base достигает уровня лучших 9-миллиардных предобученных моделей, уступая 12--20\% на RWSD и USE, 2\% на MaMuRAMu и опережая на 2\% на ruModAr.

Относительно датасета ruHHH, оценивающего способность модели определять честное, полезное и безвредное поведение (Honest, Helpful, Harmless), все модели GigaChat демонстрируют наивысшие или близкие к наивысшим результаты в своём классе: GigaChat~2~MAX, GigaChat~2~Pro и GigaChat~2 показывают лучшие или близкие к лучшим результаты по категории Harmless, незначительно уступая лидерам по Honest и Helpful ($-9\%$ до $-4\%$).

\begin{table}[H]
    \centering
    \caption{Результаты на бенчмарке MERA и ruHHH. Описания моделей доступны на странице MERA.}
    \label{tab:gc_mera}
    \resizebox{\textwidth}{!}{%
    \begin{tabular}{l|c|ccccccc}
        \toprule
        \multirow{3}{*}{\textbf{Модель}} & \multirow{3}{*}{\textbf{Итого}} & \multirow{3}{*}{\textbf{RWSD}} & \multirow{3}{*}{\textbf{ruModAr}} & \multirow{3}{*}{\textbf{USE}} & \multirow{3}{*}{\textbf{MaMuRAMu}} & \multicolumn{3}{c}{\textbf{ruHHH}} \\
        \cmidrule(lr){7-9}
         & & & & & & \textbf{Честность} & \textbf{Полезность} & \textbf{Безвредность} \\
        \midrule
        Человек (эталон) & 0.852 & 0.835 & 0.942 & 0.701 & 0.796 & 0.705 & 0.797 & 0.948 \\
        \midrule
        Claude 3.7 Sonnet & \textbf{0.682} & \textbf{0.788} & 0.919 & 0.536 & \textbf{0.89} & 0.82 & \textbf{0.864} & 0.931 \\
        \textbf{GigaChat 2 MAX} & 0.67 & 0.642 & \textbf{0.963} & \textbf{0.581} & 0.864 & 0.803 & 0.831 & \textbf{0.948} \\
        Gemini 1.5 Pro & 0.675 & 0.627 & 0.707 & 0.433 & 0.868 & 0.836 & 0.797 & 0.931 \\
        GPT-4o & 0.642 & 0.496 & 0.729 & 0.457 & 0.874 & \textbf{0.852} & 0.729 & 0.862 \\
        DeepSeek V3 & 0.677 & 0.612 & 0.718 & 0.499 & 0.882 & 0.803 & 0.763 & 0.793 \\
        \midrule
        Phi-3.5-MoE-Inst & 0.487 & 0.465 & 0.464 & 0.199 & 0.726 & 0.656 & 0.644 & 0.81 \\
        \textbf{GigaChat 2 Pro} & \textbf{0.649} & 0.665 & \textbf{0.943} & \textbf{0.534} & 0.831 & 0.803 & 0.814 & 0.897 \\
        Mixtral-8x22B-Inst & 0.486 & 0.473 & 0.523 & 0.269 & 0.747 & 0.836 & \textbf{0.881} & \textbf{0.966} \\
        Qwen2.5-72B-Inst & 0.601 & \textbf{0.715} & 0.665 & 0.32 & 0.849 & \textbf{0.869} & 0.831 & 0.897 \\
        Llama-3.1-405B-Inst & 0.59 & 0.677 & 0.573 & 0.357 & \textbf{0.868} & 0.803 & 0.864 & 0.759 \\
        \midrule
        RuadaptQwen2.5-7B & 0.536 & 0.465 & 0.492 & 0.162 & 0.751 & 0.738 & 0.78 & 0.776 \\
        \textbf{GigaChat 2} & 0.541 & 0.369 & \textbf{0.854} & \textbf{0.361} & 0.766 & 0.754 & 0.814 & \textbf{0.931} \\
        T-lite-it-1.0 & 0.552 & 0.535 & 0.493 & 0.147 & 0.775 & 0.689 & 0.797 & 0.862 \\
        \textbf{GigaChat-A3B-instruct} & 0.512 & 0.535 & 0.853 & 0.325 & 0.728 & 0.689 & 0.78 & 0.759 \\
        \textbf{GigaChat-A3B-instruct 1.5} & 0.511 & 0.512 & 0.84 & 0.32 & 0.728 & 0.689 & 0.831 & 0.793 \\
        gemma-3-27b & \textbf{0.567} & \textbf{0.588} & 0.626 & 0.328 & \textbf{0.797} & \textbf{0.82} & \textbf{0.864} & 0.914 \\
        \midrule
        gemma-2-9b & \textbf{0.453} & 0.558 & 0.592 & \textbf{0.154} & \textbf{0.689} & 0.574 & 0.627 & 0.552 \\
        \textbf{GigaChat-A3B-base} & 0.422 & 0.508 & \textbf{0.608} & 0.127 & 0.675 & 0.574 & 0.593 & 0.552 \\
        Llama-3.2-3B & 0.362 & 0.477 & 0.592 & 0.075 & 0.528 & 0.41 & 0.542 & 0.483 \\
        Yi-1.5-9B-32K & 0.428 & 0.569 & 0.516 & 0.12 & 0.516 & \textbf{0.59} & \textbf{0.661} & \textbf{0.621} \\
        Qwen1.5-7B & 0.374 & 0.558 & 0.485 & 0.056 & 0.52 & 0.541 & 0.627 & 0.603 \\
        Mistral-7B-v0.1 & 0.404 & \textbf{0.581} & 0.517 & 0.107 & 0.585 & 0.574 & 0.559 & 0.552 \\
        ruGPT-3.5 & 0.213 & 0.462 & 0.001 & 0.082 & 0.226 & 0.459 & 0.475 & 0.483 \\
        \bottomrule
    \end{tabular}%
    }
\end{table}

%% =================================================================
\section{Экологическое воздействие}\label{sec:gc_co2}
%% =================================================================

Выбросы CO$_2$ от обучения моделей вычислены по формуле Strubell et al.~\cite{strubell2019energy}:
\begin{equation}
\label{eq:co2}
    \text{CO}_2 = \frac{\text{PUE} \times \text{кВтч} \times I_{\text{CO}_2}}{1000},
\end{equation}
где PUE --- коэффициент эффективности использования энергии, кВтч --- потреблённая электроэнергия, $I_{\text{CO}_2}$ --- интенсивность выбросов углерода.

Результаты представлены в Таблице~\ref{tab:gc_co2}. 251 тысяча кг CO$_2$ приблизительно эквивалентна 157 трансатлантическим перелётам туда и обратно по маршруту Нью-Йорк --- Лондон (около 1600 кг CO$_2$ на пассажира за такой перелёт).

\begin{table}[H]
    \centering
    \caption{Выбросы CO$_2$ от обучения моделей GigaChat.}
    \label{tab:gc_co2}
    \small
    \begin{tabular}{lc}
        \toprule
        \textbf{Модель} & \textbf{CO$_2$ (кг)} \\
        \midrule
        GigaChat-A3B-base         & 251k \\
        GigaChat-A3B-instruct     & 253k \\
        GigaChat-A3B-instruct 1.5 & 255k \\
        \bottomrule
    \end{tabular}
\end{table}

%% =================================================================
\section{Выводы}\label{sec:gc_conclusions}
%% =================================================================

В настоящей главе представлено семейство LLM GigaChat --- единственная модель, разработанная с нуля на этапе предобучения специально для русского языка. Применение архитектуры MoE и специализированного токенизатора позволило создать модели, эффективно учитывающие лингвистические и культурные особенности русского языка и достигающие конкурентоспособных результатов на ведущих бенчмарках. Выпуск трёх моделей GigaChat в открытый доступ и создание удобных интерфейсов (Telegram-бот и веб-приложение для фронтирных моделей) направлены на стимулирование дальнейших исследований и промышленных приложений в области русскоязычного NLP.
